\section*{Sets and Indices}

\begin{itemize}
    \item Years relevant to decision timing: $t \in \{1,2,3\}$.
    \item Projects:
    \begin{itemize}
        \item Project $1$: short-cycle project available in each year $t=1,2,3$.
        \item Project $2$: available only at the beginning of Year 1, matures at the end of Year 2.
        \item Project $3$: available only at the beginning of Year 2, matures at the end of Year 2.
        \item Project $4$: available only at the beginning of Year 3, matures at the end of Year 3.
    \end{itemize}
\end{itemize}

\section*{Parameters}

\begin{itemize}
    \item Initial capital $C_0 = 500000$ (code: \texttt{initial\_capital}).
    \item Gross return factor of project $i$, denoted $r_i$ (code: \texttt{r1}, \texttt{r2}, \texttt{r3}, \texttt{r4}):
    \[
    r_1 = 1.20,\quad r_2 = 1.50,\quad r_3 = 1.60,\quad r_4 = 1.40.
    \]
    \item Capacity of projects with explicit upper bounds:
    \[
    U_2 = 120000,\quad U_3 = 150000,\quad U_4 = 100000
    \]
    (code: \texttt{cap2}, \texttt{cap3}, \texttt{cap4}).
    Project 1 has no explicit upper bound in the code beyond yearly cash availability.
    \item Desired liquidity reserve at the beginning of Year 3:
    \[
    R^{\min}=80000
    \]
    (code: \texttt{min\_year3\_reserve}).
    \item Penalty per yuan of reserve shortfall:
    \[
    \pi = 0.5
    \]
    (code: \texttt{reserve\_shortfall\_penalty}).
\end{itemize}

\section*{Decision Variables}

\begin{itemize}
    \item Amount invested in Project 1 at the beginning of Year $t$:
    \[
    x_{1,t} \ge 0
    \]
    with
    \[
    x_{1,1} \text{ (code: \texttt{x1\_1})},\quad
    x_{1,2} \text{ (code: \texttt{x1\_2})},\quad
    x_{1,3} \text{ (code: \texttt{x1\_3})}.
    \]
    \item Amount invested in Project 2:
    \[
    x_2 \in [0,U_2]
    \]
    (code: \texttt{x2}).
    \item Amount invested in Project 3:
    \[
    x_3 \in [0,U_3]
    \]
    (code: \texttt{x3}).
    \item Amount invested in Project 4:
    \[
    x_4 \in [0,U_4]
    \]
    (code: \texttt{x4}).
    \item Cash deliberately left uninvested at the beginning of Year 3 and carried to the end of Year 3:
    \[
    y \ge 0
    \]
    (code: \texttt{reserve\_y3}).
    \item Liquidity reserve shortfall:
    \[
    s \ge 0
    \]
    (code: \texttt{shortfall}).
\end{itemize}

\section*{Objective}

Maximize end-of-Year-3 net value after applying the liquidity shortfall penalty:
\[
\max \; r_1 x_{1,3} + r_4 x_4 + y - \pi s.
\]

Using the data values,
\[
\max \; 1.20\,x_{1,3} + 1.40\,x_4 + y - 0.5\,s.
\]

\section*{Constraints}

\paragraph{Year 1 cash availability.}
Initial capital can be allocated only to Year-1 Project 1 and Project 2:
\[
x_{1,1} + x_2 \le C_0.
\]

\paragraph{Year 2 cash availability.}
Only the proceeds from the Year-1 placement in Project 1 are available at the beginning of Year 2:
\[
x_{1,2} + x_3 \le r_1 x_{1,1}.
\]

\paragraph{Year 3 cash balance.}
At the beginning of Year 3, available cash consists of the proceeds of Year-2 Project 1 and the projects maturing at the end of Year 2; this amount is exactly split among Year-3 Project 1, Project 4, and the uninvested reserve:
\[
x_{1,3} + x_4 + y = r_1 x_{1,2} + r_2 x_2 + r_3 x_3.
\]

\paragraph{Reserve shortfall definition.}
The shortfall variable is constrained only to be at least the gap between the desired reserve and the actual carried reserve:
\[
s \ge R^{\min} - y.
\]

\paragraph{Variable bounds.}
\[
x_{1,1},x_{1,2},x_{1,3},y,s \ge 0,
\]
\[
0 \le x_2 \le U_2,\qquad 0 \le x_3 \le U_3,\qquad 0 \le x_4 \le U_4.
\]

\section*{Notes on Implementation}

\begin{itemize}
    \item The implemented model is a linear program solved by \texttt{GUROBI\_CMD} through a PuLP-compatible interface; it is not an enumeration or dynamic-programming search.
    \item The formulation follows the code exactly: Year 3 uses an equality cash-balance constraint,
    \[
    x_{1,3}+x_4+y = r_1x_{1,2}+r_2x_2+r_3x_3,
    \]
    so all cash available at the beginning of Year 3 is either invested in Year-3 opportunities or held as reserve.
    \item The shortfall condition is implemented as a one-sided inequality, not as an explicit equality:
    \[
    s \ge R^{\min}-y,\quad s\ge 0.
    \]
    Because the objective penalizes $s$, optimal solutions set $s=\max\{0,R^{\min}-y\}$, but this is an implication of optimization, not a separate model constraint.
    \item The objective counts only cash realized by the end of Year 3: returns from Project 1 placed in Year 3, returns from Project 4, and the reserved Year-3 cash carried unchanged to year-end, minus the reserve shortfall penalty.
    \item Project 1 is represented by three separate variables (\texttt{x1\_1}, \texttt{x1\_2}, \texttt{x1\_3}) rather than a single indexed variable in the code; the mathematical notation above groups them as $x_{1,t}$ for compactness.
\end{itemize}
